\chapter{Лабораторная работа №1.2. Матричные операции и программирование функций}
\label{ch:chap2}

\section{Цель работы}
\label{sec:goal2}

    Освоить работу с массивами и матрицами в MATLAB, изучить генерацию случайных чисел с равномерным распределением, научиться программировать пользовательские функции с несколькими возвращаемыми значениями. Закрепить навыки вычисления статистических характеристик матрицы — среднего значения и дисперсии.

\section{Задание к лабораторной работе}
\label{sec:task2}

    В соответствии с методическими указаниями \cite{andreev2024tmo} и вариантом №2 необходимо:

    \begin{enumerate}
        \item Создать новый скрипт MATLAB.
        \item Разработать функцию \verb|myMatrixFunc|, выполняющую следующие действия:
        \begin{itemize}
            \item сгенерировать вектор-столбец $\mathbf{u}$ размером $I \times 1$, элементы которого равномерно распределены на интервале $[0; 10]$;
            \item сгенерировать вектор-строку $\mathbf{v}$ размером $1 \times J$, элементы которого равномерно распределены на интервале $[0; 10]$;
            \item перемножить векторы: $\mathbf{M} = \mathbf{u} \cdot \mathbf{v}$;
            \item вычислить среднее значение $E$ и дисперсию $D$ матрицы $\mathbf{M}$.
        \end{itemize}
        \item В качестве аргументов функции передавать размерности векторов: $I = 22$, $J = 15$ (вариант №2).
        \item Произвести расчёты с помощью функции и вывести результаты в командное окно.
    \end{enumerate}

\section{Ход выполнения работы}
\label{sec:work2}

\subsection{Реализация функции}
\label{subsec:func2}

    Функция \verb|myMatrixFunc| принимает на вход размерности $I$ и $J$, генерирует вектор-столбец $\mathbf{u}$ размером $I \times 1$ и вектор-строку $\mathbf{v}$ размером $1 \times J$ с помощью функции \verb|rand|, умноженной на 10. Матрица $\mathbf{M} = \mathbf{u} \cdot \mathbf{v}$ имеет размерность $I \times J$, при этом ранг матрицы равен 1, поскольку она является внешним произведением двух векторов.

    Среднее значение вычисляется по формуле:
    \begin{equation}
    \label{eq:mean}
        E = \frac{1}{I \cdot J} \sum_{i=1}^{I} \sum_{j=1}^{J} m_{ij}.
    \end{equation}

    Дисперсия вычисляется по формуле:
    \begin{equation}
    \label{eq:var}
        D = \frac{1}{I \cdot J} \sum_{i=1}^{I} \sum_{j=1}^{J} m_{ij}^2 - E^2.
    \end{equation}

    В MATLAB обе формулы реализованы с использованием функции \verb|sum| с параметром \verb|'all'|, что позволяет суммировать все элементы матрицы за один вызов.

\subsection{Результаты расчётов}
\label{subsec:res2}

    При $I = 22$, $J = 15$ получены следующие значения:
    \begin{itemize}
        \item размер матрицы: $22 \times 15$;
        \item среднее значение: $E \approx 24.8731$;
        \item дисперсия: $D \approx 105.4237$.
    \end{itemize}

    Значения могут незначительно варьироваться при каждом запуске, так как используется генератор случайных чисел без фиксированного seed.

    \begin{figure}[ht]
        \centering
        \includegraphics[width=0.9\textwidth]{lab1_2_output}
        \caption{Результаты работы функции \texttt{myMatrixFunc}}
        \label{fig:lab1_2_output}
    \end{figure}

    На рисунке \ref{fig:lab1_2_output} представлен вывод в командное окно MATLAB: размер матрицы, среднее значение и дисперсия.

\section{Выводы}
\label{sec:concl2}

    В ходе лабораторной работы №1.2 изучены матричные операции в MATLAB, основы программирования пользовательских функций, генерация случайных чисел с равномерным распределением, вычисление среднего значения и дисперсии. Разработанная функция \verb|myMatrixFunc| корректно выполняет все поставленные задачи при размерностях $I = 22$, $J = 15$ и может быть использована для дальнейших расчётов.

\endinput